\textbf{Phần 1. Trắc nghiệm lựa chọn}
\setcounterpageref{ex}{0}

\begin{ex} %[1D8V33]
Tính giá trị của biểu thức: $45 + 23 - 34$.
\choice
{\True $34$}
{$33$}
{$35$}
{$32$}
\loigiai{
Thực hiện: $45 + 23 = 68$; $68 - 34 = 34$.
\textbf{Đáp số:} $34$.\par
\textbf{Đáp án của câu hỏi là A}
}
\end{ex}

\begin{ex} %[1D8V33]
Tính giá trị của biểu thức: $78 - 46 + 25$.
\choice
{$56$}
{$58$}
{\True $57$}
{$55$}
\loigiai{
Thực hiện: $78 - 46 = 32$; $32 + 25 = 57$.
\textbf{Đáp số:} $57$.\par
\textbf{Đáp án của câu hỏi là C}
}
\end{ex}

\begin{ex} %[1D8V33]
Tính giá trị của biểu thức: $99 - 54 - 23$.
\choice
{$21$}
{\True $22$}
{$23$}
{$20$}
\loigiai{
Thực hiện: $99 - 54 = 45$; $45 - 23 = 22$.
\textbf{Đáp số:} $22$.\par
\textbf{Đáp án của câu hỏi là B}
}
\end{ex}

\begin{ex} %[1D8V33]
Điền số thích hợp: $32 + \dots - 15 = 41$.
\choice
{$23$}
{$25$}
{$22$}
{\True $24$}
\loigiai{
Ta có: $32 + \dots = 41 + 15 = 56 \Rightarrow$ số cần điền là: $56 - 32 = 24$.
\textbf{Đáp số:} $24$.\par
\textbf{Đáp án của câu hỏi là D}
}
\end{ex}

\begin{ex} %[1D8V33]
Điền số thích hợp: $89 - \dots + 12 = 65$.
\choice
{\True $36$}
{$35$}
{$37$}
{$34$}
\loigiai{
Ta có: $89 - \dots = 65 - 12 = 53 \Rightarrow$ số cần điền là: $89 - 53 = 36$.
\textbf{Đáp số:} $36$.\par
\textbf{Đáp án của câu hỏi là A}
}
\end{ex}

\begin{ex} %[1D8V33]
Một người nuôi 45 con gà và 32 con vịt. Người đó bán đi 21 con gà. Hỏi người đó còn lại bao nhiêu con cả gà và vịt?
\choice
{$55$}
{$57$}
{\True $56$}
{$54$}
\loigiai{
Số gà còn lại: $45 - 21 = 24$ (con).\nTổng số con còn lại là: $24 + 32 = 56$ (con).
\textbf{Đáp số:} $56$.\par
\textbf{Đáp án của câu hỏi là C}
}
\end{ex}

\begin{ex} %[1D8V33]
Lớp 1A có 35 học sinh, lớp 1B có 34 học sinh. Có 22 học sinh tham gia đội văn nghệ. Hỏi cả hai lớp có bao nhiêu học sinh không tham gia đội văn nghệ?
\choice
{$46$}
{\True $47$}
{$48$}
{$45$}
\loigiai{
Tổng số học sinh là: $35 + 34 = 69$ (học sinh).\nSố học sinh không tham gia là: $69 - 22 = 47$ (học sinh).
\textbf{Đáp số:} $47$.\par
\textbf{Đáp án của câu hỏi là B}
}
\end{ex}

\begin{ex} %[1D8V33]
Một cửa hàng có 95 kg gạo. Buổi sáng bán 32 kg, buổi chiều bán 41 kg. Hỏi cửa hàng còn lại bao nhiêu ki-lô-gam gạo?
\choice
{$21$}
{$23$}
{$20$}
{\True $22$}
\loigiai{
Tổng số gạo bán trong ngày: $32 + 41 = 73$ (kg).\nSố gạo còn lại là: $95 - 73 = 22$ (kg).
\textbf{Đáp số:} $22$.\par
\textbf{Đáp án của câu hỏi là D}
}
\end{ex}

\begin{ex} %[1D8V33]
Tìm số tự nhiên $x$ thỏa mãn: $x + 24 = 88 - 32$.
\choice
{\True $32$}
{$31$}
{$33$}
{$30$}
\loigiai{
Vế phải: $88 - 32 = 56$. Ta có: $x + 24 = 56 \Rightarrow x = 56 - 24 = 32$.
\textbf{Đáp số:} $32$.\par
\textbf{Đáp án của câu hỏi là A}
}
\end{ex}

\begin{ex} %[1D8V33]
Tìm số tự nhiên $y$ thỏa mãn: $y - 15 = 21 + 43$.
\choice
{$78$}
{$80$}
{\True $79$}
{$77$}
\loigiai{
Vế phải: $21 + 43 = 64$. Ta có: $y - 15 = 64 \Rightarrow y = 64 + 15 = 79$.
\textbf{Đáp số:} $79$.\par
\textbf{Đáp án của câu hỏi là C}
}
\end{ex}

\begin{ex} %[1D8V33]
Một thanh sắt dài 95 cm. Bác thợ cắt bớt đi hai đoạn dài 20 cm và 35 cm. Hỏi thanh sắt còn lại dài bao nhiêu xăng-ti-mét?
\choice
{$30\text{ cm}$}
{\True $40\text{ cm}$}
{$50\text{ cm}$}
{$39\text{ cm}$}
\loigiai{
Tổng độ dài hai đoạn cắt đi là: $20 + 35 = 55$ (cm).\nĐộ dài còn lại là: $95 - 55 = 40$ (cm).
\textbf{Đáp số:} $40$.\par
\textbf{Đáp án của câu hỏi là B}
}
\end{ex}

\begin{ex} %[1D8V33]
Tính tổng của số lớn nhất có hai chữ số và hiệu của 55 với 55.
\choice
{$98$}
{$100$}
{$97$}
{\True $99$}
\loigiai{
Hiệu của 55 và 55 là 0. Số lớn nhất có hai chữ số là 99: $99 + 0 = 99$.
\textbf{Đáp số:} $99$.\par
\textbf{Đáp án của câu hỏi là D}
}
\end{ex}

\textbf{Phần 2. Trắc nghiệm đúng sai}
\setcounterpageref{ex}{0}

\begin{ex} %[1D8V33]
Xét tính đúng hoặc sai của các khẳng định toán học sau đây:
\choiceTF
{\True Kết quả của bài toán «Tìm một số biết rằng lấy số đó cộng với 33 rồi trừ đi 21 thì được 56» là $44$.}
{\True Kết quả của bài toán «Tìm một số biết rằng lấy 98 trừ đi số đó rồi cộng thêm 11 thì được 65» là $44$.}
{Bình có 45 viên bi, An có 23 viên bi. Bình cho Cúc 12 viên bi. lúc này số bi của Bình nhiều hơn số bi của An $11$ viên bi.}
{Hai số có tổng bằng 79, số bé là 34. số lớn là số $44$ .}
\loigiai{
\begin{itemize}
    \item Ý a) Đúng: Số cần tìm là: $56 + 21 - 33 = 44$.
\textbf{Đáp số:} $44$.
    \item Ý b) Đúng: Ta có: $98 - x = 65 - 11 = 54 \Rightarrow x = 98 - 54 = 44$.
\textbf{Đáp số:} $44$.
    \item Ý c) Sai: Kết quả đúng là $10$, không phải $11$. Số bi của Bình còn lại là: $45 - 12 = 33$ (viên bi).\nSố bi của Bình nhiều hơn của An là: $33 - 23 = 10$ (viên bi).
\textbf{Đáp số:} $10$.
    \item Ý d) Sai: Kết quả đúng là $45$, không phải $44$. Số lớn là: $79 - 34 = 45$.
\textbf{Đáp số:} $45$.
\end{itemize}
\begin{itemize}
    \item a) \textbf{Đúng}.
    \item b) \textbf{Đúng}.
    \item c) \textbf{Sai}.
    \item d) \textbf{Sai}.
\end{itemize}
}
\end{ex}

\begin{ex} %[1D8V33]
Xét tính đúng hoặc sai của các bài toán và phép tính sau đây:
\choiceTF
{Hai số có hiệu bằng 23, số lớn là 86. số bé là số $62$ .}
{Giá trị của phép tính $11 + 22 + 33 + 11$ là $76$.}
{\True Giá trị của phép tính $99 - 11 - 22 - 33$ là $33$.}
{\True Bình có 45 viên bi, An có 23 viên bi. Bình cho Cúc 12 viên bi. lúc này số bi của Bình nhiều hơn số bi của An $10$ viên bi.}
\loigiai{
\begin{itemize}
    \item Ý a) Sai: Kết quả đúng là $63$, không phải $62$. Số bé là: $86 - 23 = 63$.
\textbf{Đáp số:} $63$.
    \item Ý b) Sai: Kết quả đúng là $77$, không phải $76$. Thực hiện: $11 + 22 = 33$; $33 + 33 = 66$; $66 + 11 = 77$.
\textbf{Đáp số:} $77$.
    \item Ý c) Đúng: Thực hiện: $99 - 11 = 88$; $88 - 22 = 66$; $66 - 33 = 33$.
\textbf{Đáp số:} $33$.
    \item Ý d) Đúng: Số bi của Bình còn lại là: $45 - 12 = 33$ (viên bi).\nSố bi của Bình nhiều hơn của An là: $33 - 23 = 10$ (viên bi).
\textbf{Đáp số:} $10$.
\end{itemize}
\begin{itemize}
    \item a) \textbf{Sai}.
    \item b) \textbf{Sai}.
    \item c) \textbf{Đúng}.
    \item d) \textbf{Đúng}.
\end{itemize}
}
\end{ex}

\textbf{Phần 3. Trả lời ngắn}
\setcounterpageref{ex}{0}

\begin{ex} %[1D8V33]
Một trang trại nuôi 68 con vịt và 31 con ngan. Sau đó trang trại mua thêm 10 con ngan nữa. Hỏi đàn gia cầm có tất cả bao nhiêu con ngan?
\shortans{41}
\loigiai{
\textbf{Đáp số:} $41$\par
Số con ngan có tất cả là: $31 + 10 = 41$ (con ngan).
\textbf{Đáp số:} $41$.
}
\end{ex}

\begin{ex} %[1D8V33]
Số con vịt nhiều hơn số con ngan lúc này là bao nhiêu con?
\shortans{27}
\loigiai{
\textbf{Đáp số:} $27$\par
Số vịt nhiều hơn số ngan là: $68 - 41 = 27$ (con).
\textbf{Đáp số:} $27$.
}
\end{ex}

\begin{ex} %[1D8V33]
Tìm số thích hợp điền vào dấu chấm: $55 + 23 = 99 - \dots$.
\shortans{21}
\loigiai{
\textbf{Đáp số:} $21$\par
Vế trái: $55 + 23 = 78$. Số cần điền: $99 - 78 = 21$.
\textbf{Đáp số:} $21$.
}
\end{ex}

\begin{ex} %[1D8V33]
Một cây vải dài 85 mét. Buổi sáng bán 30 mét, buổi chiều bán 25 mét. Cây vải còn lại bao nhiêu mét?
\shortans{30}
\loigiai{
\textbf{Đáp số:} $30$\par
Tổng số vải đã bán là: $30 + 25 = 55$ (m).\nSố mét vải còn lại là: $85 - 55 = 30$ (mét).
\textbf{Đáp số:} $30$.
}
\end{ex}

\begin{ex} %[1D8V33]
Tính giá trị biểu thức: $34 + 43 - 21 - 12$.
\shortans{44}
\loigiai{
\textbf{Đáp số:} $44$\par
Thực hiện: $34 + 43 = 77$; $77 - 21 = 56$; $56 - 12 = 44$.
\textbf{Đáp số:} $44$.
}
\end{ex}

\begin{ex} %[1D8V33]
Tìm số tự nhiên có hai chữ số giống nhau biết số đó trừ đi 22 thì được 44.
\shortans{66}
\loigiai{
\textbf{Đáp số:} $66$\par
Số đó là: $44 + 22 = 66$.
\textbf{Đáp số:} $66$.
}
\end{ex}
